% !TeX root = ../../main.tex

Das Problem mit großen Systemen ist, dass diese immer diverser werden.
Es gibt viele unterschiedliche Komponenten, die in einem großen System laufen sollen.
Die einzelnen Komponenten sind eine Datenbank, SAP und vieles mehr.
Dies alles in einem System ist zu komplex, um vernünftig integriert, optimiert und gewartet zu werden.
Deshalb hat sich die Forschung mit autonomen Systemen auseinandergesetzt.

Ein autonomes System wird durch mehrere Konzepte definiert.
Bei \emph{Self-configuration} sollen sich Komponenten und Systeme anhand \emph{high-level policies}\footnote{Richtlinien
 der Firma, die festlegen, was ein System können muss.} automatisch konfigurieren.
\emph{Self-optimization} besagt, dass Komponenten und Systeme selbst nach Möglichkeiten suchen sich zu verbessern.
Außerdem soll sich ein System von selbst gegen Angriffe schützen.
Dies fällt unter den Bereich \emph{Self-protection}.
Der letzte Punkt ist \emph{Self-healing}, welches beschreibt, dass ein System automatisiert Software und Hardware
Probleme erkennt und von selbst repariert.
Retry-Mechanismen sind ein Teil des \emph{Self-healing} und sorgen dafür, dass ein System autonomer agiert~\cite{Kephart2003}.
